\section{Mode A6 -- Mise en conditions initiales}
\label{secModeA6}

\subsection{Conditions d'entrée}

Le mode A6 est atteint depuis le mode A5, une fois la remise en route
autorisée par l'opérateur, ou depuis le mode F4 (mode manuel), lorsque
l'opérateur demande le retour aux conditions initiales de l'installation
(à condition qu'aucun mouvement ne soit alors en cours en mode F4).

\subsection{Comportement attendu}

L'objectif de ce mode est d'amener le robot dans sa position de référence,
au-dessus du guichet, sans jamais entraver la circulation des palettes sur
le convoyeur pendant que le robot se déplace vers cette position.

\begin{itemize}
    \item \textbf{Robot} : le robot doit rejoindre sa position de
    référence au-dessus du guichet, en restant sous puissance pendant tout
    le déplacement. Le trajet emprunté ne doit à aucun moment risquer
    d'entrer en collision avec une palette présente sur le convoyeur.
    \item \textbf{Convoyeur / guichet} : l'accès au guichet doit rester
    interdit pendant toute la durée de la mise en conditions initiales.
    \item \textbf{Gestion d'un arrêt demandé par l'opérateur} : si un
    arrêt est demandé par l'opérateur pendant que le robot est en
    mouvement, le robot doit être immobilisé (tout en restant sous
    puissance) et doit reprendre sa progression vers la position de
    référence dès que cette demande d'arrêt n'est plus active.
    \item \textbf{IHM} : la mise en conditions initiales en cours doit être
    signalée à l'opérateur par une signalisation dédiée. En cas d'arrêt
    demandé pendant ce mode, cette situation particulière doit également
    être signalée distinctement, ainsi que le bouton concerné.
\end{itemize}

\subsection{Conditions de sortie}

\begin{itemize}
    \item Vers le mode \textbf{A1} (\emph{Arrêt en conditions initiales}),
    une fois que le robot a atteint sa position de référence au-dessus du
    guichet.
\end{itemize}
